% =====================================================================
% Section I: Introduction
% =====================================================================
\section{Introduction}
\label{sec:intro}

\IEEEPARstart{L}{owering} the radiation dose of a computed tomography (CT)
examination lowers the number of detected photons, and the reconstruction
problem becomes both noisier and more ill conditioned. Statistical iterative
reconstruction addresses this by modelling the photon statistics explicitly,
starting from maximum likelihood expectation maximisation (MLEM)
\cite{shepp1982}, and by adding a prior that restrains the noise the
likelihood alone cannot control. Total variation \cite{sidky2008},
anisotropic diffusion \cite{perona1990}, plug and play priors
\cite{venkatakrishnan2013} and fractional order penalties
\cite{zhang2016fractional} all follow this pattern, and all share the same
practical burden: at least one strength parameter must be set, and the
appropriate value depends on the dose, the anatomy and the acquisition
geometry.

Learned reconstruction removes that burden by moving the prior into the
training data. Post-processing networks \cite{chen2017redcnn,yang2018wgan},
unrolled physics informed architectures
\cite{adler2018lpd,wu2021drone,xiang2021fista}, architecture searched
denoisers \cite{lu2023m3nas} and, more recently, diffusion priors used as
generative regularisers \cite{chung2023dps,gao2024corediff,liu2024mar} now
define the state of the art on the public low-dose benchmarks
\cite{moen2021mayo,leuschner2021lodopab}. Their accuracy in the setting they
were trained for is not in question. What is in question is what happens
outside it. Deep reconstruction networks have been shown to be unstable under
small perturbations and to miss structures that were not represented in
training \cite{antun2020,genzel2023}, and a controlled benchmark of low-dose
denoisers found that newer architectures do not consistently outperform older
ones once training and evaluation are held fixed \cite{eulig2024}. For a
modality whose acquisitions vary across scanners, protocols and patient
populations, a method whose behaviour changes with the training distribution
carries a risk that accuracy on one dataset does not describe.

This work takes the opposite route. We keep the reconstruction free of
training data, and we remove the tuning burden from the classical side by
making the regulariser select itself from the measurements. The starting
point is our earlier fractional-order notch filter \cite{singh2026spl}, in
which a spectral operator combines a fractional-order low-pass envelope with
narrowband stop bands placed at frequencies where structured artefacts
concentrate. That formulation required the fractional order and the notch
positions to be supplied. Here both are inferred from the data at hand: the
order is chosen by a Poisson discrepancy criterion on an unbounded ladder of
candidates, and the notches are detected by a robust spectral rule with an
acquisition band guard, so that a single frozen algorithm runs unchanged
across doses, anatomies and geometries.

The resulting method is evaluated against a jury that spans the field:
analytic reconstruction, four classical iterative priors, a training free
neural prior \cite{ulyanov2018dip,baguer2020dip}, two supervised networks
\cite{chen2017redcnn,adler2018lpd} and a diffusion posterior sampler
\cite{chung2023dps}, on two public datasets, two geometries, four dose levels
and ten noise realisations per configuration, with paired statistics at two
levels of aggregation. The evaluation reports what the numbers say, including
where the proposed method is not the best.

\subsection*{Contributions}

\begin{enumerate}
\item \textbf{A reconstruction with no run time parameters.} The fractional
order is selected by a Poisson discrepancy criterion evaluated on a geometric
ladder that is not bounded a priori, the notch set is detected by a robust
rule that combines a sectoral background model, a spatial coherence test and
an acquisition band guard, and the iteration itself is a damped
multiplicative step followed by a nonexpansive spectral step. Nothing is
tuned per dataset, per dose or per geometry.

\item \textbf{Stability that the theory predicts and the data confirm.} The
spectral operator is self-adjoint and nonexpansive for every admissible
order, so the regularised iteration cannot amplify. Empirically, across 200
photon starved reconstructions per method the proposed iteration produced no
catastrophic failure, whereas the nonlinear spatial priors produced between
three and sixteen, and an unregularised MLEM run at a fixed budget diverged
in fan beam geometry while the same engine with the spectral step stayed
stable.

\item \textbf{Behaviour under distribution shift.} Transferred without any
adjustment to a second public dataset, the proposed method reproduces its own
accuracy to within a few hundredths of a decibel, while the supervised and
generative baselines lose between 2.7 and 7.0 dB and the ranking of the
methods inverts. The size of the degradation orders with the amount of
physics each architecture contains.

\item \textbf{A statistically rigorous and released benchmark.} Ten methods,
two datasets, two geometries, four dose levels, ten realisations, paired
Wilcoxon tests at the realisation and the slice level, a task based
assessment with a channelised Hotelling observer, and a released code base
with an automated verification suite that checks the operators, the observer
and the reconstruction pipeline against their reference definitions.
\end{enumerate}

The paper is organised as follows. Section \ref{sec:related} reviews related
work. Section \ref{sec:method} develops the method. Section
\ref{sec:results} reports the experiments and discusses them, including the
regimes in which classical priors remain preferable. Section
\ref{sec:conclusion} concludes.
